% Hafta 3 -- alıştırma gövdesi (ders notu ve ayrı alıştırma kağıdı tarafından \input edilir)
\section{Çözümlü Alıştırmalar}
Her soruyu önce kendiniz deneyin, sonra çözümü okuyun.

\begin{alistirma}
\(\ab{6a3}\) üç basamaklı sayısı \(11\) ile bölünebiliyorsa \(a\) kaçtır?
\end{alistirma}
\begin{cozum}
Birler basamağından başlayarak \(+,-,+\): \ \(3-a+6=9-a\). Bu sayı \(11\)'in katı olmalıdır.\\
\(0\le a\le 9\) olduğundan \(0\le 9-a\le 9\); bu aralıktaki tek \(11\) katı \(0\)'dır: \(9-a=0\Rightarrow a=9\) \ (\(693=11\cdot 63\)).
\end{cozum}

\begin{alistirma}
\(\ab{2a5b}\) dört basamaklı sayısı \(5\) ile ve \(9\) ile bölünebiliyorsa \(a\) en çok kaçtır?
\end{alistirma}
\begin{cozum}
\(5\) ile bölünme: \(b\in\{0,5\}\). \quad \(9\) ile bölünme: rakamlar toplamı \(2+a+5+b=7+a+b\) sayısı \(9\)'un katı.\\
\(b=0\): \(7+a\in\{9\}\) (\(7\le 7+a\le 16\)) \(\Rightarrow a=2\) \quad (\(2250\))\\
\(b=5\): \(12+a\in\{18\}\) (\(12\le 12+a\le 21\)) \(\Rightarrow a=6\) \quad (\(2655\))\\
\(a\) en çok \(6\)'dır. \ Kontrol: \(2655=5\cdot 531=9\cdot 295\).
\end{cozum}

\begin{alistirma}
\(\ab{3a4b}\) dört basamaklı sayısı \(36\) ile bölünebiliyorsa bu koşulu sağlayan kaç farklı sayı vardır?
\end{alistirma}
\begin{cozum}
\(36=4\cdot 9\) (aralarında asal çarpanlar; \(36=6\cdot 6\) veya \(3\cdot 12\) kullanılamaz).\\
\(4\) ile bölünme: \(\ab{4b}\in\{40,44,48\}\Rightarrow b\in\{0,4,8\}\).\\
\(9\) ile bölünme: \(3+a+4+b=7+a+b\) sayısı \(9\)'un katı.\\
\(b=0\): \(7+a=9\Rightarrow a=2\) \ (\(3240\)); \quad \(b=4\): \(11+a=18\Rightarrow a=7\) \ (\(3744\)); \quad \(b=8\): \(15+a=18\Rightarrow a=3\) \ (\(3348\)).\\
Cevap: \(3\) farklı sayı. \ (Kontrol: \(3240=36\cdot 90\), \(3348=36\cdot 93\), \(3744=36\cdot 104\).)
\end{cozum}

\begin{alistirma}
\(A=2^3\cdot 3^2\cdot 5\) ve \(B=2\cdot 3^3\cdot 5^2\) ise \(\text{EBOB}(A,B)\) ve \(\text{EKOK}(A,B)\) kaçtır?
\end{alistirma}
\begin{cozum}
EBOB: ortak asal çarpanlar en küçük üsleriyle: \ \(2^{1}\cdot 3^{2}\cdot 5^{1}=90\).\\
EKOK: bütün asal çarpanlar en büyük üsleriyle: \ \(2^{3}\cdot 3^{3}\cdot 5^{2}=8\cdot 27\cdot 25=5400\).\\
Kontrol: \(A=360\), \(B=1350\); \ \(\text{EBOB}\cdot\text{EKOK}=90\cdot 5400=486\,000=360\cdot 1350=A\cdot B\) \checkmark
\end{cozum}

\begin{alistirma}
\(a\) ve \(b\) aralarında asal iki pozitif tam sayıdır ve \(a\cdot b=90\)'dır. \(a+b\) en az kaçtır?
\end{alistirma}
\begin{cozum}
\(90=2\cdot 3^2\cdot 5\). \(a\) ile \(b\) aralarında asal olduğundan her asal çarpan bütün kuvvetiyle \textbf{yalnız birinde} bulunur; \(3^2=9\) bölünemez. Parçalar: \(2\), \(9\), \(5\).\\
Olası ikililer: \(\{1,90\}\to 91\), \ \(\{2,45\}\to 47\), \ \(\{5,18\}\to 23\), \ \(\{9,10\}\to 19\).\\
\(a+b\) en az \(9+10=19\). \ (\(\{6,15\}\) çarpımı \(90\) olsa da \(\text{EBOB}(6,15)=3\) olduğundan uygun değildir.)
\end{cozum}

\begin{alistirma}
Kenarları \(96\) m ve \(120\) m olan dikdörtgen bir bahçenin çevresine, her köşede birer ağaç olacak ve ardışık ağaçlar arası uzaklıklar eşit olacak biçimde ağaç dikilecektir. En az kaç ağaç gerekir?
\end{alistirma}
\begin{cozum}
En az ağaç için aralık en büyük olmalı ve her kenarı tam bölmeli: \(96=2^5\cdot 3\), \(120=2^3\cdot 3\cdot 5\) \(\Rightarrow\) \(\text{EBOB}(96,120)=2^3\cdot 3=24\) m.\\
Çevre: \(2\cdot(96+120)=432\) m. \ Kapalı çevre: ağaç sayısı \(=432\div 24=18\) (\(+1\) eklenmez).\\
Kenar kenar kontrol: \(96\div 24=4\), \(120\div 24=5\) aralık \(\Rightarrow\) \(2\cdot(4+5)=18\) aralık \(=18\) ağaç.
\end{cozum}

\begin{alistirma}
Üç otobüs hattından otobüsler sırasıyla \(20\), \(30\) ve \(45\) dakikada bir kalkmaktadır. Üç hattan saat 07{:}15'te birlikte otobüs kalktıktan sonra en erken saat kaçta yeniden birlikte kalkar?
\end{alistirma}
\begin{cozum}
Birlikte kalkış aralığı üç sürenin en küçük ortak katıdır:
\(20=2^2\cdot 5\), \(30=2\cdot 3\cdot 5\), \(45=3^2\cdot 5\) \(\Rightarrow\) \(\text{EKOK}=2^2\cdot 3^2\cdot 5=180\) dakika \(=3\) saat; \ 07{:}15\(\,+\,3\) saat \(=\) 10{:}15.
\end{cozum}

\begin{alistirma}
\(4\), \(6\) ve \(9\) ile bölündüğünde sırasıyla \(2\), \(4\) ve \(7\) kalanını veren üç basamaklı en küçük doğal sayı kaçtır?
\end{alistirma}
\begin{cozum}
Her kalan bölenden \(2\) eksiktir: \(4-2=2\), \(6-2=4\), \(9-2=7\). O halde \(x+2\) sayısı \(4\), \(6\) ve \(9\)'un ortak katıdır.\\
\(\text{EKOK}(4,6,9)=2^2\cdot 3^2=36\) \(\Rightarrow\) \(x=36k-2\): \ \(34,\ 70,\ 106,\dots\)\\
Üç basamaklı en küçük: \(x=106\). \ Kontrol: \(106=4\cdot 26+2=6\cdot 17+4=9\cdot 11+7\).
\end{cozum}

\kendinizoncesi
\section{Alıştırmalar (Kendiniz Çözün)}

\begin{enumerate}[label=\textbf{\arabic*.},leftmargin=*,itemsep=4pt]
\item \(\ab{24a}\) üç basamaklı sayısı \(4\) ile bölünebiliyorsa \(a\) hangi değerleri alabilir?
\item \(\ab{5a8}\) üç basamaklı sayısı \(3\) ile bölünebiliyorsa \(a\)'nın alabileceği değerlerin toplamı kaçtır?
\item \(9\,876\,543\) sayısının \(9\) ile ve \(5\) ile bölümünden kalanlar kaçtır?
\item \(3\,125\,764\) sayısının \(8\) ile bölümünden kalan kaçtır?
\item \(\text{EBOB}(72,108)\) ve \(\text{EKOK}(72,108)\) kaçtır?
\item İki pozitif tam sayının EBOB'u \(9\), EKOK'u \(270\)'tir. Sayılardan biri \(54\) ise diğeri kaçtır?
\item \(\ab{a4b}\) biçimindeki üç basamaklı sayılardan \(6\) ile bölünebilen en büyüğü kaçtır?
\item Boyutları \(150\) cm \(\times\) \(210\) cm olan dikdörtgen bir zemin, eş kare fayanslarla boşluksuz ve fayans kesilmeden döşenecektir. En az kaç fayans gerekir?
\item Kenarları \(15\) cm ve \(20\) cm olan eş dikdörtgen kartonlar, aynı yönde ve boşluksuz yan yana dizilerek bir kare oluşturulacaktır. En az kaç karton gerekir?
\item \(\ab{a1b}\) biçimindeki, \(45\) ile bölünebilen üç basamaklı sayıların toplamı kaçtır?
\item \(a<b\) olmak üzere EBOB'u \(4\), EKOK'u \(60\) olan kaç farklı \((a,b)\) pozitif tam sayı ikilisi vardır?
\item Bir kutudaki kalemler \(8\)'er, \(12\)'şer ve \(15\)'er sayıldığında her seferinde \(3\) kalem artmaktadır. Kalem sayısı \(200\)'den az olduğuna göre kutuda en çok kaç kalem vardır?
\end{enumerate}

\subsection*{Cevap Anahtarı}
{\small
\begin{multicols}{2}
\begin{enumerate}[label=\textbf{\arabic*.},leftmargin=*,itemsep=1pt]
\item \(\ab{4a}\in\{40,44,48\}\Rightarrow a\in\{0,4,8\}\).
\item \(13+a\) \(3\)'ün katı \(\Rightarrow a\in\{2,5,8\}\); toplam \(15\).
\item Rakam toplamı \(42\): \(9\) ile kalan \(6\); \(5\) ile kalan \(3\).
\item Son üç basamak \(764=8\cdot 95+4\Rightarrow\) kalan \(4\).
\item \(72=2^3\cdot 3^2\), \(108=2^2\cdot 3^3\); EBOB \(36\), EKOK \(216\).
\item \(54\cdot b=9\cdot 270\Rightarrow b=45\).
\item \(a=9\); \(b\) çift, \(13+b\) \(3\)'ün katı; en büyük \(948\).
\item \(\text{EBOB}(150,210)=30\); \(5\cdot 7=35\) fayans.
\item \(\text{EKOK}(15,20)=60\); \(4\cdot 3=12\) karton.
\item \(b=0\Rightarrow 810\); \(b=5\Rightarrow 315\); toplam \(1125\).
\item \(a=4x\), \(b=4y\), \(xy=15\): \((4,60)\), \((12,20)\); \(2\) ikili.
\item \(x-3\), \(120\)'nin katı: \(3,\ 123\); en çok \(123\).
\end{enumerate}
\end{multicols}}
